\documentclass[11pt]{article}
\usepackage[utf8]{inputenc}
\usepackage[T1]{fontenc}
\usepackage{amsmath,amssymb,amsthm}
\usepackage[margin=2.5cm]{geometry}
\usepackage[numbers]{natbib}
\usepackage{hyperref}

\newtheorem{theorem}{Theorem}
\newtheorem{lemma}[theorem]{Lemma}
\newtheorem{proposition}[theorem]{Proposition}
\theoremstyle{definition}
\newtheorem{definition}[theorem]{Definition}
\newtheorem{assumption}{Assumption}

\title{Early Stopping for Random Planted Forests}
\author{Oskar Allerslev \and Markus Baltzer}
\date{\today}

\begin{document}
\maketitle

\begin{abstract}
To be written.
\end{abstract}

\section{Introduction}
Random planted forests \citep{hiabu2023rpf} are interpretable tree ensembles whose
fitted function decomposes into low-order interaction terms. Early stopping rules for
regression trees based on the residual norm were developed by \citet{miftachov2025early}.
This paper studies early stopping for random planted forests.

\section{Setting}

\section{Early stopping rule}

\section{Main results}

\section{Simulations}

\bibliographystyle{plainnat}
\bibliography{../literature/references}

\end{document}
